\documentclass{article}
\usepackage[T1]{fontenc}
\usepackage{lmodern}
\usepackage{graphicx}
\usepackage{amsmath,amssymb}
\usepackage[a4paper,margin=2cm]{geometry}
\usepackage[absolute,overlay]{textpos}
\setlength{\TPHorizModule}{1mm}
\setlength{\TPVertModule}{1mm}
\usepackage[indonesian]{babel}
\usepackage{hyperref}
\setlength{\emergencystretch}{2em}
% unit: Halaman 23

\begin{document}

% === Halaman 23 (llm) ===

\section*{Elliptic Curve Discrete Logarithm Problem (ECDLP)}

\begin{itemize}
    \item Menghitung $kP = Q$ mudah, tetapi menghitung $k$ jika diketahui $P$ dan $Q$ adalah sulit. Inilah ECDLP yang menjadi dasar ECC.
    \item ECDLP dirumuskan sebagai berikut:
    
    \textcolor{red}{Diberikan $P$ dan $Q$ adalah dua buah titik di kurva eliptik, carilah integer $k$ sedemikian sehingga $Q = kP$}
    
    \item Secara komputasi sulit menemukan $k$, jika $k$ adalah bilangan yang besar. Nilai $k$ adalah logaritma diskrit dari $Q$ dengan basis $P$. \textcolor{red}{*)}
    \item Pada algoritma ECC, $Q$ adalah kunci publik, $k$ adalah kunci privat, dan $P$ sembarang titik pada kurva eliptik.
\end{itemize}

\centering
\textcolor{red}{Catatan: ingatlah $kP = P^k$, sehingga $Q = kP = P^k$, $k$ adalah logaritma diskrit dari $Q$}

\end{document}
